\documentclass{situationist-preprint}

\title{Divided Selves}
\msid{PID-15-a}
\msfamily{PID-15}
\msdiscipline{Personal Identity and the Self}
\mstype{Principal manuscript}
\mskeywords{dissociative identity disorder; unity of the self; partial unity; alters; co-consciousness; amnesia}

\begin{document}
\maketitle

\begin{abstract}
Can a self be divided? Dissociative identity disorder, in which a person presents two or more distinct identity states with gaps in memory between them, has been read as a case of several selves in one body and, by sceptics, as a single self enacting roles. We argue for a third reading: dissociative conditions involve the partial unity of a single self. Extending MND-15's account of partial unity of consciousness, we call a self's states partially unified when co-accessibility, availability together to one act of self-ascription or self-regulation, links any two of them by a chain but does not hold between every two. Shared skills, one-way amnesias, co-conscious alters and the experimentally observed transfer of information between identity states show that the states are so linked, and also that co-accessibility fails between some of them. A self so structured is one self, partially unified.
\end{abstract}

\section{The problem}

The current diagnostic criteria describe dissociative identity disorder as a disruption of identity marked by two or more distinct personality states, with recurrent gaps in the recall of everyday events and personal information \citep{APA2013}. The identity states, or alters, differ in name, manner and self-description, and some report no memory of what others did. Three readings have been proposed. On the first, each alter is a self, and the condition is a plurality of selves in one body. On the second, associated with the sociocognitive model, there is one self enacting several roles under the influence of therapeutic and cultural expectations \citep{Spanos1994}. On the third, which \citet{Braude1995} defends, there is a single subject underlying the alters, which are dissociated systems within it.

We defend a version of the third reading and make it precise using the notion of partial unity that MND-15-a develops for consciousness in split-brain cases.

\section{Framework}

\begin{definition}[Co-accessibility]\label{d:co}
Two states of a subject, experiential, mnemonic or practical, are \emph{co-accessible} if both are available together to a single act of self-ascription or self-regulation: the subject can ascribe both to herself in one judgment, or bring both to bear in the control of one action.
\end{definition}

\begin{definition}[Unity of a self]\label{d:unity}
A self's states are \emph{unified} if every two of them are co-accessible. They are \emph{partially unified} if every two are linked by a chain of co-accessible pairs, but not every two are co-accessible. A \emph{self}, in the counting sense, is a maximal set of states linked by such chains.
\end{definition}

Definition~\ref{d:unity} follows the structure of MND-15-a's definition of the unity of consciousness, with co-accessibility in the place of co-consciousness. Partial unity is the failure of transitivity: states $s_1$ and $s_3$ may each be co-accessible with $s_2$ without being co-accessible with each other. Two selves require two sets of states with no chain between them. The self here is the construction of PID-04, whose unity can be a matter of degree as a simple substance's could not.

\begin{principle}[Evidence of co-accessibility]\label{p:ev}
If information acquired in one identity state influences judgment or action in another, or if a state is reported, from within one identity state, as one's own and as accompanying a state of another, the states concerned are linked by co-accessibility. If information acquired in one state is unavailable in another under conditions that would reveal it if it were available, they are not co-accessible.
\end{principle}

\section{The argument}

\begin{lemma}\label{l:linked}
In dissociative identity disorder, the states of the different identity states are linked by chains of co-accessibility.
\end{lemma}

\begin{proof}
Four kinds of evidence satisfy the first clause of Principle~\ref{p:ev}. Alters share skills, language and general knowledge, which each brings to bear in its own actions. Amnesia between alters is frequently one-way: one alter reports awareness of what another does while the second reports no awareness of the first, so the first alter's self-ascriptions take in states of both. Some alters report being co-conscious with others, observing the actions of the alter in executive control as they occur. And experimental studies of patients whose alters report mutual amnesia have found that information learned by one alter influences the recall and recognition performance of another \citep{Huntjens2003}. Hence any two identity states are linked, directly or through a third, by co-accessible pairs.
\end{proof}

\begin{lemma}\label{l:fail}
In dissociative identity disorder, co-accessibility fails between some states.
\end{lemma}

\begin{proof}
Some autobiographical episodes undergone by one alter are not self-ascribed by another under conditions, such as direct questioning in therapy without motive to conceal, that would reveal them if they were available; and alters maintain distinct self-representations that are not combined in a single act of self-ascription. By the second clause of Principle~\ref{p:ev}, these states are not co-accessible.
\end{proof}

\begin{theorem}\label{t:main}
Dissociative conditions involve the partial unity of a single self: a self can be divided, but only partially.
\end{theorem}

\begin{proof}
By Lemma~\ref{l:linked}, all the states belong to one maximal set linked by chains of co-accessibility, and so, by Definition~\ref{d:unity}, to one self. By Lemma~\ref{l:fail}, not every two of them are co-accessible. By Definition~\ref{d:unity}, the self's states are partially unified.
\end{proof}

The result is the self-level counterpart of MND-15's conclusion for consciousness. It also confirms the verdict, recorded in PID-03-a, that the division found in a single person is at most partial, and contrasts it with fission, in which two persons arise and no chain of co-accessibility links them after the operation.

\section{Objections and replies}

\begin{objection}[Many selves]
Alters have distinct names, histories and self-conceptions. If a self is a structure of self-representation, each alter is one, and there are several selves.
\end{objection}

\begin{reply}
Each alter is a self-representation, and some are elaborate. But by Definition~\ref{d:unity} selves are counted by chains of co-accessibility, not by self-representations, and Lemma~\ref{l:linked} shows that the alters are so linked. The objection counts the descriptions a self gives of itself rather than the self that gives them.
\end{reply}

\begin{objection}[Enactment]
On the sociocognitive model, alters are roles enacted in response to cues. There is no division to explain, only a performance \citep{Spanos1994}.
\end{objection}

\begin{reply}
The objection concerns how alters arise, not how the resulting states are structured. If alters are socially shaped constructions, so is the self as a whole (PID-04), and being constructed does not make a thing unreal. The question is whether the failures of co-accessibility in Lemma~\ref{l:fail} are real, and the experimental evidence of selective failures of access supports that they are, even while showing that they are less complete than patients report. That is exactly what partial unity predicts and full division does not.
\end{reply}

\begin{objection}[Full unity]
The evidence of transfer between alters shows that the self is fully unified and that the amnesias are reporting failures.
\end{objection}

\begin{reply}
Transfer shows that chains of co-accessibility exist, not that every pair is co-accessible. Implicit influence on performance is compatible with the unavailability of the same information to self-ascription, and the latter is what Lemma~\ref{l:fail} requires. Full unity would need every two states to be available to one act of self-ascription, which the evidence does not show.
\end{reply}

\section{Conclusion}

Dissociative identity disorder is neither a plurality of selves nor a single self in disguise. Its identity states are bound into one self by chains of co-accessibility and kept apart by failures of access between some of them: the relation that unifies them is not transitive, and that is what a partially unified self is.

\begin{thebibliography}{99}\small
\bibitem[American Psychiatric Association(2013)]{APA2013} American Psychiatric Association (2013). \emph{Diagnostic and Statistical Manual of Mental Disorders}, 5th edn. American Psychiatric Publishing.
\bibitem[Braude(1995)]{Braude1995} Braude, S.~E. (1995). \emph{First Person Plural: Multiple Personality and the Philosophy of Mind}, rev. edn. Rowman \& Littlefield.
\bibitem[Huntjens et~al.(2003)]{Huntjens2003} Huntjens, R.~J.~C., Postma, A., Peters, M.~L., Woertman, L., and van der Hart, O. (2003). Interidentity amnesia for neutral, episodic information in dissociative identity disorder. \emph{Journal of Abnormal Psychology}, 112(2), 290--297.
\bibitem[Spanos(1994)]{Spanos1994} Spanos, N.~P. (1994). Multiple identity enactments and multiple personality disorder: A sociocognitive perspective. \emph{Psychological Bulletin}, 116(1), 143--165.
\end{thebibliography}

\end{document}
